\documentclass{amsart}
% For dealing with references we use the comment environment
\usepackage{verbatim}
\newenvironment{reference}{\comment}{\endcomment}
%\newenvironment{reference}{}{}
\newenvironment{slogan}{\comment}{\endcomment}
\newenvironment{history}{\comment}{\endcomment}

% For commutative diagrams we use Xy-pic
\usepackage[all]{xy}

% We use 2cell for 2-commutative diagrams.
\xyoption{2cell}
\UseAllTwocells

% We use multicol for the list of chapters between chapters
\usepackage{multicol}

% This is generally recommended for better output
\usepackage{lmodern}
\usepackage[T1]{fontenc}

% For cross-file-references

% Package for hypertext links:
\usepackage{hyperref}

% For any local file, say "hello.tex" you want to link to please

% Theorem environments.
%
\theoremstyle{plain}
\newtheorem{theorem}[subsection]{Theorem}
\newtheorem{proposition}[subsection]{Proposition}
\newtheorem{lemma}[subsection]{Lemma}

\theoremstyle{definition}
\newtheorem{definition}[subsection]{Definition}
\newtheorem{example}[subsection]{Example}
\newtheorem{exercise}[subsection]{Exercise}
\newtheorem{situation}[subsection]{Situation}

\theoremstyle{remark}
\newtheorem{remark}[subsection]{Remark}
\newtheorem{remarks}[subsection]{Remarks}

\numberwithin{equation}{subsection}

% Macros
%
\def\lim{\mathop{\mathrm{lim}}\nolimits}
\def\colim{\mathop{\mathrm{colim}}\nolimits}
\def\Spec{\mathop{\mathrm{Spec}}}
\def\Hom{\mathop{\mathrm{Hom}}\nolimits}
\def\Ext{\mathop{\mathrm{Ext}}\nolimits}
\def\SheafHom{\mathop{\mathcal{H}\!\mathit{om}}\nolimits}
\def\SheafExt{\mathop{\mathcal{E}\!\mathit{xt}}\nolimits}
\def\Sch{\mathit{Sch}}
\def\Mor{\mathop{\mathrm{Mor}}\nolimits}
\def\Ob{\mathop{\mathrm{Ob}}\nolimits}
\def\Sh{\mathop{\mathit{Sh}}\nolimits}
\def\NL{\mathop{N\!L}\nolimits}
\def\CH{\mathop{\mathrm{CH}}\nolimits}
\def\proetale{{pro\text{-}\acute{e}tale}}
\def\etale{{\acute{e}tale}}
\def\QCoh{\mathit{QCoh}}
\def\Ker{\mathop{\mathrm{Ker}}}
\def\Im{\mathop{\mathrm{Im}}}
\def\Coker{\mathop{\mathrm{Coker}}}
\def\Coim{\mathop{\mathrm{Coim}}}

% Boxtimes
%
\DeclareMathSymbol{\boxtimes}{\mathbin}{AMSa}{"02}

%
% Macros for moduli stacks/spaces
%
\def\QCohstack{\mathcal{QC}\!\mathit{oh}}
\def\Cohstack{\mathcal{C}\!\mathit{oh}}
\def\Spacesstack{\mathcal{S}\!\mathit{paces}}
\def\Quotfunctor{\mathrm{Quot}}
\def\Hilbfunctor{\mathrm{Hilb}}
\def\Curvesstack{\mathcal{C}\!\mathit{urves}}
\def\Polarizedstack{\mathcal{P}\!\mathit{olarized}}
\def\Complexesstack{\mathcal{C}\!\mathit{omplexes}}
% \Pic is the operator that assigns to X its picard group, usage \Pic(X)
% \Picardstack_{X/B} denotes the Picard stack of X over B
% \Picardfunctor_{X/B} denotes the Picard functor of X over B
\def\Pic{\mathop{\mathrm{Pic}}\nolimits}
\def\Picardstack{\mathcal{P}\!\mathit{ic}}
\def\Picardfunctor{\mathrm{Pic}}
\def\Deformationcategory{\mathcal{D}\!\mathit{ef}}

\setlength{\emergencystretch}{3em}
\begin{document}
\title[Stacks Project, Tag 0368]{629 --- The Jacobson property is fppf local}
\maketitle
\noindent Copyright (C) 2005--2025 Johan de Jong. Stacks Project authors. Source: \href{https://stacks.math.columbia.edu/tag/0368}{Tag 0368}.

\noindent Editorial difficulty note (not author text). Difficulty H1: The substantive descent argument first passes to a quasi-finite faithfully flat refinement, then combines going down with the Jacobson criterion to force an infinite fibre contradiction. This geometric descent proof goes beyond routine ring-theoretic qualification checks..

\noindent Editorial provenance note (not author text). History: excerpt from revision \texttt{a0444\allowbreak 6e57e\allowbreak c1fbc\allowbreak 252a8\allowbreak 71afc\allowbreak ec775\allowbreak 2fb28\allowbreak 07b14}, descent.tex, lines 4784--4850. Reference presentation was adapted; original-author context and core support blocks were retained or added. The mathematical wording is unchanged. Licensed under GNU FDL 1.2 or later; no invariant sections or cover texts. See ../licenses/COPYING. Context blocks reproduce author definitions and supporting results; they are not additional counted problems.

\begin{lemma}
\label{lemma-Jacobson-local-fppf}
The property $\mathcal{P}(S) =$``$S$ is Jacobson'' is local
in the fppf topology.
\end{lemma}

\begin{proof}
We will use Lemma \href{https://stacks.math.columbia.edu/tag/0349}{lemma descending properties (Tag 0349)}.
First we note that ``being Jacobson'' is local
in the Zariski topology. This is
Properties, Lemma \href{https://stacks.math.columbia.edu/tag/01P4}{lemma locally jacobson (Tag 01P4)}.
Next, we show that if $S' \to S$ is a flat, finitely presented
morphism of affines and $S$ is Jacobson, then $S'$ is
Jacobson. This is
Morphisms, Lemma \href{https://stacks.math.columbia.edu/tag/02J5}{lemma Jacobson universally Jacobson (Tag 02J5)}.
Finally, we have to show that if $f : S' \to S$ is a surjective
flat, finitely presented morphism of affines and $S'$ is
Jacobson, then $S$ is Jacobson. Say $S = \Spec(A)$ and
$S' = \Spec(B)$ and $S' \to S$ given by $A \to B$.
Then $A \to B$ is finitely presented and faithfully flat.
Moreover, the ring $B$ is Jacobson, see
Properties, Lemma \href{https://stacks.math.columbia.edu/tag/01P4}{lemma locally jacobson (Tag 01P4)}.

\medskip\noindent
By Algebra, Lemma \href{https://stacks.math.columbia.edu/tag/034Z}{lemma fppf fpqf (Tag 034Z)} there exists a diagram
$$
\xymatrix{
B \ar[rr] & & B' \\
& A \ar[ru] \ar[lu] &
}
$$
with $A \to B'$ finitely presented, faithfully flat and quasi-finite.
In particular, $B \to B'$ is finite type, and we see from
Algebra, Proposition \href{https://stacks.math.columbia.edu/tag/00GB}{proposition Jacobson permanence (Tag 00GB)}
that $B'$ is Jacobson. Hence we may assume that $A \to B$ is quasi-finite
as well as faithfully flat and of finite presentation.

\medskip\noindent
Assume $A$ is not Jacobson to get a contradiction.
According to Algebra, Lemma \href{https://stacks.math.columbia.edu/tag/034J}{lemma characterize jacobson (Tag 034J)}
there exists a nonmaximal prime $\mathfrak p \subset A$ and
an element $f \in A$, $f \not \in \mathfrak p$ such that
$V(\mathfrak p) \cap D(f) = \{\mathfrak p\}$.

\medskip\noindent
This leads to a contradiction as follows. First let
$\mathfrak p \subset \mathfrak m$ be a maximal ideal of $A$.
Pick a prime $\mathfrak m' \subset B$ lying over $\mathfrak m$
(exists because $A \to B$ is faithfully flat, see
Algebra, Lemma \href{https://stacks.math.columbia.edu/tag/00HQ}{lemma ff rings (Tag 00HQ)}).
As $A \to B$ is flat, by going down see
Algebra, Lemma \href{https://stacks.math.columbia.edu/tag/00HS}{lemma flat going down (Tag 00HS)},
we can find a prime $\mathfrak q \subset \mathfrak m'$ lying over
$\mathfrak p$. In particular we see that $\mathfrak q$ is not
maximal. Hence according to
Algebra, Lemma \href{https://stacks.math.columbia.edu/tag/034J}{lemma characterize jacobson (Tag 034J)} again
the set $V(\mathfrak q) \cap D(f)$ is infinite
(here we finally use that $B$ is Jacobson).
All points of $V(\mathfrak q) \cap D(f)$ map to
$V(\mathfrak p) \cap D(f) = \{\mathfrak p\}$. Hence the
fibre over $\mathfrak p$ is infinite. This contradicts the
fact that $A \to B$ is quasi-finite (see
Algebra, Lemma \href{https://stacks.math.columbia.edu/tag/00PM}{lemma quasi finite (Tag 00PM)}
or more explicitly
Morphisms, Lemma \href{https://stacks.math.columbia.edu/tag/02NH}{lemma quasi finite (Tag 02NH)}).
Thus the lemma is proved.
\end{proof}
\end{document}
